\chapter{YÊU CẦU HỆ THỐNG}\label{chap:yeucau}

Yêu cầu được chia thành ba nhóm: chức năng (hệ thống phải làm gì), phi chức năng (làm tốt đến mức nào) và môi
trường sử dụng (cần gì để vận hành). Mỗi yêu cầu chức năng ứng với một module có thể chạy độc lập.

\section{Yêu cầu chức năng}

\begin{enumerate}[label=\textbf{FR\arabic*.}, leftmargin=3.2em]
\item \textbf{Thu thập:} lấy metadata và PDF từ arXiv API theo keyword, chuyên mục, khoảng ngày hoặc mã bài;
  không tải trùng.
\item \textbf{Nạp tài liệu:} nhận PDF, trích tiêu đề, tác giả, tóm tắt, mã arXiv và danh sách tài liệu tham khảo
  có mã arXiv.
\item \textbf{Kiểm tra trùng:} theo mã băm SHA-256 của file, mã arXiv, tiêu đề chuẩn hóa và tiêu đề gần giống.
\item \textbf{Embedding và lưu trữ:} chia đoạn, sinh vector, lưu vào vector DB và cập nhật đồ thị một cách tự
  động sau khi nạp.
\item \textbf{Tìm kiếm:} tìm kiếm lai BM25 + vector, gộp bằng Reciprocal Rank Fusion, lọc theo năm.
\item \textbf{Graph RAG:} mở rộng kết quả qua đồ thị (tác giả, chuyên mục, khái niệm, trích dẫn, tương đồng).
\item \textbf{Chatbot:} hai chế độ Chat và Research, trả lời có trích dẫn, ghi nhớ hội thoại, hỏi theo bài vừa
  thêm.
\item \textbf{Giao diện:} chat streaming, lịch sử hội thoại, kéo-thả PDF, danh sách nguồn và trực quan hóa đồ thị
  của từng câu trả lời.
\end{enumerate}

\section{Yêu cầu phi chức năng}

\begin{table}[H]
\centering\small\setstretch{1.15}
\begin{tabularx}{\linewidth}{@{}P{0.18\linewidth}LL@{}}
\toprule
\textbf{Nhóm} & \textbf{Yêu cầu} & \textbf{Cách đáp ứng}\\
\midrule
Hiệu năng & Tìm kiếm dưới 1 giây với kho 10.000 bài; nạp một PDF dưới 30 giây trên CPU & Chỉ mục HNSW (ChromaDB), BM25 trong bộ nhớ, embedding theo lô\\
Khả năng mở rộng & Embedding hàng chục nghìn bài & Chạy theo lô trên GPU Kaggle, xuất thư mục dữ liệu để dùng lại\\
Tính module & Mỗi module chạy độc lập & Gói Python \code{scholarlens.*}, mỗi module có lệnh riêng\\
Tính đúng đắn & Không bịa nguồn & Prompt buộc trích dẫn [n] từ ngữ cảnh; chế độ không LLM trả về trích đoạn\\
Tính linh hoạt & Đổi LLM hoặc model embedding & Cấu hình qua biến môi trường hoặc giao diện: API tương thích OpenAI (vLLM, Gemini, Ollama), Anthropic, hoặc không LLM\\
Tính khả dụng & Vẫn dùng được khi mất kết nối LLM & Tự quay về chế độ trích đoạn; tìm kiếm và đồ thị chạy cục bộ\\
Tính bảo trì & Dữ liệu tái lập được & Vector store và đồ thị dựng lại được toàn bộ từ SQLite\\
\bottomrule
\end{tabularx}
\caption{Yêu cầu phi chức năng}\label{tab:nfr}
\end{table}

\section{Yêu cầu môi trường sử dụng}

Máy demo không cần GPU. Mọi tính toán nặng nằm trên máy GPU, và máy GPU có thể là một phiên Kaggle miễn phí.

\begin{table}[H]
\centering\small\setstretch{1.15}
\begin{tabularx}{\linewidth}{@{}P{0.18\linewidth}LL@{}}
\toprule
\textbf{Thành phần} & \textbf{Yêu cầu tối thiểu} & \textbf{Cấu hình đã dùng khi phát triển}\\
\midrule
Máy GPU (LLM + embedding hàng loạt) & GPU NVIDIA có compute capability $\geq$ 7.5; tổng VRAM $\geq$ 16 GB cho model 7B nén AWQ ($\geq$ 8 GB cho model 3B); Linux, Python 3.10+, vLLM, có Internet & Kaggle Notebook, 2 $\times$ Tesla T4 (2 $\times$ 16 GB), Python 3.13, vLLM 0.32.0\\
Máy demo (web app) & CPU 4 nhân, RAM 8 GB, khoảng 3 GB ổ đĩa; Docker, hoặc Python 3.10+ & Windows 10, Python 3.13, chỉ CPU\\
Trình duyệt & Chrome, Edge hoặc Firefox bản hiện hành & Microsoft Edge\\
Mạng & Máy demo gọi ra ngoài qua HTTPS (tới tunnel trycloudflare, và arXiv API khi thêm PDF); không cần mở cổng vào & Mạng gia đình thông thường\\
Tài khoản & Kaggle đã xác minh để bật GPU và Internet; Hugging Face token nếu dùng model có giới hạn truy cập & Kaggle; Qwen2.5 không cần token\\
\bottomrule
\end{tabularx}
\caption{Yêu cầu môi trường sử dụng}\label{tab:moitruong}
\end{table}

Hai ràng buộc phải giữ đồng nhất giữa hai máy: cùng model embedding (\code{intfloat/multilingual-e5-base}) và cùng
phiên bản \code{chromadb} (1.5.9), để dữ liệu tạo trên máy GPU đọc được trên máy demo. Khi không có máy GPU, hệ
thống vẫn chạy ở chế độ trích đoạn: tìm kiếm, đồ thị và thêm PDF hoạt động bình thường, chỉ thiếu câu trả lời do
LLM viết.
